% ==================================================================
% AUDITING FRAMEWORK (one page, shared by all audit documents)
% The validity chain (benchmark -> measurement -> evaluation ->
% claim), the validity aspects attached to each step, and the audit
% dimensions that operationalise them. Claim-level validities are
% explicitly out of scope. Eval-independent: edit here and every
% audit document picks up the change.
% ==================================================================
\phantomsection
\addcontentsline{toc}{section}{Auditing Framework}
\label{sec:framework}

\begin{adjustwidth}{-10mm}{0mm}

% --- Title row -----------------------------------------------------
\noindent
{\LARGE\itshape \textbf{Auditing Framework}}\hfill
{\small\itshape\color{black!55}version \aauditframework}

\smallskip

% --- Intro ----------------------------------------------------------
\noindent
\begin{dsbox}[top=1mm, bottom=1mm, before skip=3pt, after skip=3pt]{Validity aspects}
\footnotesize
An evaluation supports a chain of increasingly ambitious inferences:
running a \textbf{benchmark} produces a \textbf{measurement}, the
measurement is interpreted as an \textbf{evaluation} of some
capability, and the evaluation is enlisted to support \textbf{claims}.
Each link can fail in its own way, so validity threats attach at
different points of the chain. The audit dimensions on the next page
operationalise the validity aspects below the claim line; each
dimension is assessed in the Audit Findings. The validity aspects
associated with claims (criterion, external, and consequential
validity) depend on the specific claim being made and its deployment
context, not on the eval alone -- they are \textbf{beyond the scope of
this audit}.
\end{dsbox}

\smallskip

% --- Diagram ---------------------------------------------------------
% Chips name audit dimensions (scorecard row labels). "Task
% specification" and "Tasks and datasets" are both assessed under the
% Tasks-and-datasets dimension, which straddles the measurement and
% evaluation rungs.
\newcommand{\fwchip}[1]{{\setlength{\fboxsep}{2.5pt}\setlength{\fboxrule}{0.5pt}%
  \fcolorbox{black!40}{white}{\footnotesize #1}}}
\newcommand{\fwchipg}[1]{{\setlength{\fboxsep}{2.5pt}\setlength{\fboxrule}{0.5pt}%
  \fcolorbox{black!35}{black!4}{\footnotesize\color{black!55}#1}}}
\noindent
\begin{tikzpicture}[
    aspectbox/.style={rectangle, rounded corners=1.5mm, draw=brandTeal,
      line width=0.6pt, fill=brandTeal!4, inner sep=2.5mm,
      text width=98mm, align=left, anchor=north west},
    greybox/.style={rectangle, rounded corners=1.5mm, draw=black!45,
      line width=0.6pt, fill=black!3, inner sep=2.5mm,
      text width=98mm, align=left, anchor=north west},
    rung/.style={rectangle, rounded corners=1.5mm, fill=brandTeal,
      text=white, minimum width=46mm, minimum height=13mm,
      text width=42mm, align=center},
    rungx/.style={rung, fill=black!45},
    tolink/.style={->, line width=0.7pt, black!60},
    ladder/.style={->, line width=1pt, black!60}]

  % --- claim tier (out of scope) -------------------------------------
  \node[greybox] (claimbox) at (0mm,0mm) {%
    {\small\bfseries\color{black!60}Claim validities -- beyond audit scope}\\[3pt]
    \fwchipg{Criterion validity}\ \fwchipg{External validity}\
    \fwchipg{Consequential validity}\\[3pt]
    {\scriptsize\itshape\color{black!55}Assessed against a specific
     claim and deployment context, not the eval alone; explicitly out
     of scope for this audit.}};
  \node[rungx] (claim) at (139mm,-13mm) {%
    {\small\bfseries Claim}\\[1pt]
    {\scriptsize\itshape what the eval is used to argue}};

  % --- evaluation tier ------------------------------------------------
  \node[aspectbox] (construct) at (0mm,-36mm) {%
    {\small\bfseries\color{brandTeal}Construct validity}\\[3pt]
    \fwchip{{\bfseries\color{brandTeal}A}~Construct validity}\ \fwchip{{\bfseries\color{brandTeal}I}~Informativeness}};
  \node[aspectbox] (content) at (0mm,-56mm) {%
    {\small\bfseries\color{brandTeal}Content validity}\\[3pt]
    \fwchip{{\bfseries\color{brandTeal}B}~Content validity}};
  \node[rung] (evaluation) at (139mm,-53mm) {%
    {\small\bfseries Evaluation}\\[1pt]
    {\scriptsize\itshape the interpretation of the score}};

  % --- measurement tier -----------------------------------------------
  \node[aspectbox] (impl) at (0mm,-78mm) {%
    {\small\bfseries\color{brandTeal}Implementation validity}\\[3pt]
    \fwchip{{\bfseries\color{brandTeal}C}~Task Specification}\ \fwchip{{\bfseries\color{brandTeal}D}~Scaffolding}\ \fwchip{{\bfseries\color{brandTeal}E}~Harness}\\[3pt]
    \fwchip{{\bfseries\color{brandTeal}F}~Environment}\ \fwchip{{\bfseries\color{brandTeal}G}~Grading}\ \fwchip{{\bfseries\color{brandTeal}H}~Resource limits}};
  \node[rung] (measurement) at (139mm,-89mm) {%
    {\small\bfseries Measurement}\\[1pt]
    {\scriptsize\itshape the score produced by a run}};

  % --- benchmark ------------------------------------------------------
  \node[rung] (benchmark) at (139mm,-118mm) {%
    {\small\bfseries Benchmark}\\[1pt]
    {\scriptsize\itshape tasks, data, harness, grading}};

  % --- ladder ----------------------------------------------------------
  \draw[ladder] (benchmark) -- (measurement);
  \draw[ladder] (measurement) -- (evaluation);
  \draw[ladder] (evaluation) -- (claim);

  % --- aspect -> rung ---------------------------------------------------
  \draw[tolink, dashed] (claimbox.east) -- (claim.west);
  \draw[tolink] (construct.east) -- (evaluation.west);
  \draw[tolink] (content.east) -- (evaluation.west);
  \draw[tolink] (impl.east) -- (measurement.west);
\end{tikzpicture}

\end{adjustwidth}
